\documentclass[11pt,letterpaper]{article}

\usepackage[margin=1in]{geometry}
\usepackage{times}
\usepackage{amsmath,amssymb}
\usepackage{setspace}
\usepackage[hidelinks]{hyperref}
\usepackage{microtype}

\pagestyle{empty}
\setstretch{1.08}
\setlength{\parskip}{0.55em}
\setlength{\parindent}{0pt}

\newcommand{\abssection}[1]{\textbf{#1.}~}

\begin{document}

\begin{center}
{\LARGE\bfseries Structural Influence:}\\[0.35em]
{\large\bfseries Measuring the Hidden Wiring of an NBA Offense}\\[1em]
{\large Rami Zheman}\\[0.25em]
{\normalsize\textit{Independent Sports Analytics Researcher}\\
Alexandria, VA}
\end{center}

\vspace{0.6em}

\abssection{Introduction}
At the 2024 trade deadline, Dallas gave up a first-round pick for
Daniel Gafford, a backup center averaging eleven points and eight
rebounds, and reached the Finals with him starting at center.
Existing analytics, built around production, efficiency, and usage,
cannot explain what Dallas was actually buying. NBA front offices
routinely pay for players whose value is organizational rather than
statistical, glue guys who hold a system together, yet there has been
no quantitative way to identify that role or measure how completely
an offense's identity depends on it. This paper asks whether
offensive organization can be measured directly, whether some
low-usage players are load-bearing in a way box scores cannot see,
and whether removing them changes what a team actually is, not just
how well it scores.

\abssection{Methods}
Each team-season is represented as a network connecting players to
the play types they initiate, built from public NBA play-by-play data
across three seasons, 2022--23 through 2024--25, covering roughly
650{,}000 possessions and 924 rotation player-seasons. A team's
structural fingerprint is a standardized matrix capturing how often
the same players jointly initiate different play-type pairs. A
player's structural influence is the change in that fingerprint when
his initiations are removed, isolated from usage to reveal residual
influence. Crossing usage and residual influence identifies shapers,
low-usage players whose absence reorganizes an offense as much as a
star's. This measure is then validated against real basketball:
actual on-court and off-court minutes are compared directly, testing
whether observed team identity shifts when a player sits.

\abssection{Results}
Shapers, players like Gafford and Rudy Gobert, carry structural
influence comparable to primary organizers despite scoring less than
half as much. When a shaper actually leaves the floor, the offense's
real structure shifts far more than when a usage-matched non-shaper
leaves, a larger identity change in 81 percent of matched pairs,
confirmed with real on-off data rather than a hypothetical. The
hypothetical removal measure and the real on-off shift are strongly
correlated: a number computed by deleting a player from a dataset
reliably predicts something that happens on an actual court. Shapers
reshape what a team runs, more cuts and pick-and-rolls, fewer
spot-ups and drives, and part of this influence travels with them
after a trade. None of this moves the scoreboard: structural
completeness does not raise scoring efficiency. A team can look
entirely different with a shaper on the floor and score the same
either way.

\abssection{Conclusion}
An NBA offense can change its entire operating identity around one
low-usage player without the scoreboard ever registering it. That
distinction, between what makes a team score and what makes a team
itself, has been invisible to conventional analytics and is now
measurable. For the industry, this offers a reproducible way to
identify which players a roster cannot afford to lose in any sense
other than points, informing how front offices value trades, how
coaches plan around injuries, and how teams protect the identity a
system was built around.

\end{document}
